% One-time imported draft. Editable native source; no Markdown regeneration.
\section{Render a full frame before
publishing}\label{009-sorting-visualizer-p04-render-a-full-frame-before-publishing}

The frame is persistent memory. A smaller bar after an exchange needs
black stores to erase the old top, and an inactive column must be
cleared even if it once contained a bar. Rendering only the white cells
would leave a history of earlier states in the image.

The row-major frame cursor advances four bytes for every cell. A
separate column index selects the input word, and a row index decides
whether that word\textquotesingle s bar reaches the current cell. The
value is loaded only for an active column. Every path then reaches the
same word store with either white or black.

All sixty-four words are ready before the public label. The stores
themselves are not an atomic frame transaction, so a running IDE can
repaint a partial render. Pausing at \texttt{sort\_ready} gives the
complete state time to appear. The future capture process must verify
both that model boundary and the visible panel; a fixed sleep alone is
insufficient.
